\documentclass[10pt,a4paper]{amsart}
\usepackage[margin=19mm]{geometry}
\usepackage{amsmath,amssymb,mathtools}
\usepackage[scr=rsfs,cal=euler]{mathalfa}
\usepackage{xfrac}
\usepackage[unicode,hidelinks,bookmarks=false]{hyperref}
\newcommand{\ie}{i.e.\ }
\newcommand{\cf}{cf.\ }
\newcommand{\resp}{respectively\ }
\newcommand{\Subsection}{\subsection}
\providecommand{\nobreakdash}{\nobreak\hbox{-}}
\setlength{\emergencystretch}{2em}
\multlinegap0pt
\usepackage{bm}
\usepackage{bbm}
\usepackage{dsfont}
\usepackage{xspace}
\usepackage{cancel}
\usepackage{mathtools}
\usepackage{amsthm}
\makeatletter
\@ifundefined{Theorem}{\newtheorem{Theorem}{Theorem}}{}
\@ifundefined{Proposition}{\newtheorem{Proposition}[Theorem]{Proposition}}{}
\makeatother
\title[Higher-Order Linear Differential Equations for Unitary Matri]{Higher-Order Linear Differential Equations for Unitary Matrix Integrals: Applications and Generalisations}
\author{Peter J. Forrester, Fei Wei (with an appendix by Folkmar Bornemann)}
\date{Source: arXiv:2508.20797v3 (CC BY 4.0)}
\begin{document}
\maketitle

\noindent\emph{Source and licence.} Higher-Order Linear Differential Equations for Unitary Matrix Integrals: Applications and Generalisations. Version arXiv:2508.20797v3 (CC BY 4.0), distributed under the Creative Commons Attribution 4.0 International licence, as recorded in the arXiv licence metadata for this exact version and shown on the abstract page https://arxiv.org/abs/2508.20797v3. Full rights notice: licenses/arxiv-2508.20797-CC-BY-4.0.txt.\par\smallskip

\noindent\textit{Editorial scope.} Counted unit: the Proposition whose source label is \texttt{P3.3} in the article (source0.tex lines 481--495, source label \texttt{P3.3}), delivered together with the article's own complete original proof of it (source0.tex lines 497--540, one \texttt{proof} environment of 44 lines). No mathematics is written, repaired or re-derived: statement and proof are the article's own text line for line. Required context retained uncounted: 6.1a (lines 328--332); 6.1b (lines 374--377). Every definition, display and label of those lines is retained unchanged. Omitted from this package and not redistributed: 1--327, 333--373, 378--480, 496--496, 541--1411. Nothing inside a retained excerpt is omitted or altered: every retained body is delivered exactly as the article's lines are written, so no omission receipt of any kind accompanies this package. Citations: the retained excerpts carry no citation command at all, so no bibliography entry of the article is transcribed and no reference is redistributed. Reference adaptation: none was needed and none was made; every reference inside the delivered unit resolves inside it, so no unresolved reference or citation is produced. Printed slips kept literal: no printed word, symbol, hypothesis or number of the counted statement or of its proof was altered. Layout and class: the article's own class is \texttt{sigma} with packages including ifpdf, mathtools, bbm, braket, stmaryrd, makecell; the wrapper is the lane's pinned \texttt{amsart} editorial wrapper at 10pt on A4, which restates the theorem-like environments the retained text needs and adds the article's own macro definitions that the retained text needs (none), with extra wrapper packages (bm, bbm, dsfont, xspace, cancel, mathtools). The delivered numbering is the unit's own. Assets: no retained span contains a figure environment, a graphics-inclusion command, a PDF-page inclusion, a table or a TikZ picture; no image file is copied into this package, the package declares no asset dependency and no image file is redistributed with it. Licence: version arXiv:2508.20797v3 of this article is distributed under the Creative Commons Attribution 4.0 International licence, as recorded in the arXiv licence metadata for this exact version and shown on the abstract page https://arxiv.org/abs/2508.20797v3. A full rights notice is delivered at licenses/arxiv-2508.20797-CC-BY-4.0.txt. Characters: every retained excerpt is pure ASCII, so the delivered engine, XeLaTeX with its default Latin Modern text font, reports no missing character.
\begin{gather}
(n - p) x H_{p+1}(x)
= ( (n - p) x + B_p ) H_p(x) + x {{\rm d} \over {\rm d}x} H_p(x) - D_p H_{p-1}(x), \nonumber\\
 p=0,\dots, n,\label{6.1a}
\end{gather}

\begin{equation}\label{6.1b}
 (n - p) x I_{p+1}(x)
= B_p I_p(x) + x {{\rm d} \over {\rm d}x} I_p(x) - p I_{p-1}(x), \qquad p=0,\dots, n,
\end{equation}

\begin{Proposition}\label{P3.3}
 Let $I_n(x)$ be specified via the differential-difference system
 \eqref{6.1b}. For some polynomials $\{ q_m(x) \}_{m=0,1,\dots,n-1}$, we have
 \begin{gather}
\left ( x^n {{\rm d}^{n+1} \over {\rm d} x^{n+1}} + q_n(x) x^{n-1} {{\rm d}^{n} \over {\rm d} x^{n}} + q_{n-1}(x) x^{n-2}{{\rm d}^{n-1} \over{\rm d} x^{n-1}}+ \cdots +
 q_2(x) x{{\rm d}^2 \over{\rm d} x^2}+
 q_1(x) {{\rm d} \over{\rm d} x} + q_0(x) \right )\nonumber\\
 \qquad\times
 I_n(x) =0.\label{MD3a2}
 \end{gather}
 The degree in $x$ of the coefficient of each
 ${{\rm d}^{j} \over{\rm d} x^{j}}$ is $[(n-1+j)/2]$, thus restricting the degrees of the~$q_m(x)$. Furthermore, when~$\{ B_j \}$
 {in \eqref{6.1b}}
 are all integers, the coefficients of these polynomials are also integers.
 \end{Proposition}

 \begin{proof}
 Let\footnote{This notation is not to be confused with $D_p$ in~\eqref{6.1a}.} $D_x := x {{\rm d} \over{\rm d} x}$.
 From \eqref{6.1b} in the case $p=n$, we read off that
 \begin{equation}\label{3.8a}
 n I_{n-1}(x) = (B_n + D_x) I_{n}(x).
 \end{equation}
 Using this in the case $p=n-1$ of \eqref{6.1b} then gives
 \begin{equation}\label{3.8b}
 n ( n - 1 ) I_{n-2}(x) = ((B_{n-1} + D_x)
 (B_{n} + D_x) - n x )
 I_{n}(x).
 \end{equation}
 Using both \eqref{3.8a} and \eqref{3.8b} in the case $p=n-2$ of
\eqref{6.1b} then gives
 \begin{gather}
 n ( n - 1 )(n-2) I_{n-3}(x)\nonumber\\
 \qquad = ( (B_{n-2} + D_x) ((B_{n-1} + D_x)
 (B_{n} + D_x) - n x )
 -2 x (n-1) (B_{n} + D_x) )
 I_{n}(x).\label{3.8c}
 \end{gather}
 We see that the differential operators on the right-hand sides of \eqref{3.8a}, \eqref{3.8b}, \eqref{3.8c} have the structure
 \begin{gather*}%\label{3.8d}
 x {{\rm d} \over {\rm d}x} + c_0,\qquad
 x^2 {{\rm d}^2 \over {\rm d}x^2} + d_0 x {{\rm d} \over {\rm d}x}+
 (c_1 + c_2 x),\\
 x^3 {{\rm d}^3 \over {\rm d}x^3} + e_0 x^2 {{\rm d}^2 \over {\rm d}x^2}+
 (d_1 + d_2 x) x {{\rm d} \over {\rm d}x} +
 (c_3 + c_4x),
\end{gather*}
{for certain coefficients.}

Repeating this procedure down to the case $p=1$ of \eqref{6.1b} allows each $I_j(x)$ ($j=0,\dots,n)$ to be written as a differential operator of degree $n-j$ acting on $I_n(x)$. In particular,
$n! I_1(x)$ has the structure of a differential operator of the form \eqref{MD3a2} with $n \mapsto n - 1$ and an extra factor of~$x$ multiplying each of the terms involving a derivative. Similarly, $n! I_0(x)$ has the structure of a~differential operator of the form in
\eqref{MD3a2} as written, except for
an extra factor of $x$ multiplying each of the terms involving a derivative.

The case $p=0$ of \eqref{6.1b} gives
 \begin{equation}\label{op3a}
{{\rm d} \over{\rm d} x} I_0(x) - n I_1(x) = 0.
\end{equation}
From the established formulas for $I_0(x)$, $I_1(x)$, this implies \eqref{MD3a2}. That the coefficients of the polynomials $\{ q_j(x) \}$ are integers for $\{ B_j\}$ integers is immediate from the structure of the right-hand sides of
\eqref{3.8a}--\eqref{3.8c} as extends to $I_0(x)$, $I_1(x)$.
\end{proof}

\end{document}
